\chapter{紫外発散と超対称性による解決}

\lectureoverview{超対称性の形式的な数学を導入する前に、それが解決しようとしている問題について正確に述べておくべきである。ここでは、位置空間と運動量空間の両方で紫外発散を記述し、次元解析によってスカラーおよびフェルミオンの伝播関数の短距離挙動を導き、スカラー質量への主要なループ補正を見積もる。スカラーループとフェルミオンループは同じ二次のカットオフ依存性を持つが、符号は反対である。その相殺が偶発的ではなく構造的となるために、どのような原理がそれらの結合、多重度、質量を十分に一致させうるのかを問うとき、超対称性は説得力を持つ。}

\section{なぜ発散に戻るのか？}

当初は超対称性の仕組みを直ちに始めるつもりであった。しかし、抽象化を1回の講義だけ遅らせ、動機を再構築する方がよい。以前の発散と質量の繰り込みについての説明は急ぎ足で進んだため、何を修正する必要があるのかを正確に見ておかないと、超対称性は恣意的な数学的発明のように思われるであろう。

ファインマン図は、密接に関連する2つの仕方で素粒子物理学に現れる。最もよく知られているのは、散乱振幅を組織化するというものである。入射粒子と出射粒子を指定し、必要な次数のすべての図が表す振幅を足し合わせる。また、\emph{有効ラグランジアン}の構成を組織化するというものでもある。ここで``有効''という語は、異常に成功しているという褒め言葉ではない。その記述が近似的であること、すなわち、アクセスできない短距離の詳細が消去または積分除去され、その影響が残ったパラメータと相互作用に符号化されていることを意味する。

そのような記述にはそれぞれ有効範囲がある。その範囲は、ここでは\(\delta\)で表される最短の信頼できる距離、または最大の信頼できる運動量
\begin{equation}
\Lambda\sim \delta^{-1}.
\end{equation}
によって指定できる。この2つのカットオフは別々の仮定ではない。同じ紫外極限をフーリエ双対的に記述したものである。

\section{1つの伝播関数、2つの記述}

伝播関数とは、ある時空点で粒子を生成または挿入し、別の時空点でそれを除去または検出する振幅である。2つの点を\(x^\mu\)と\(y^\mu\)とし、その隔たりを
\begin{equation}
\Delta^\mu=y^\mu-x^\mu.
\end{equation}
と書く。小文字の記号\(\delta\)はカットオフのために取っておくので、隔たりと正則化を混同することはない。

並進不変性により、自由伝播関数は\(x\)と\(y\)に\(\Delta\)を通してのみ依存する。同じ対象はフーリエ変換できる：
\begin{equation}
G(\Delta)
=\int\frac{d^4k}{(2\pi)^4}\,e^{-ik\cdot\Delta}\,\widetilde G(k).
\end{equation}

「位置にフーリエ共役なものは何か？」という問い``が場に投げかけられると、答えは運動量である''。したがって、ループが始点に戻る伝播関数であるという位置空間の記述は、ループの周りを循環できるすべての運動量にわたる積分であるという運動量空間の記述と等価である。運動量空間の規則では各頂点で運動量保存を課す必要があるが、位置空間では現在のパワーカウントの議論がより透明に保たれる。

この等価性は、何かを計算する前に紫外問題を同定する。位置空間では、ループは一致点極限\(\Delta\to0\)を探る。運動量空間では、極限\(|k|\to\infty\)を探る。大きな運動量は短い波長を意味するので、短距離特異性と高運動量発散は、異なる座標における同じ現象である。

この講義では最も単純なスピンクラスを扱う：スカラー場、フェルミオン場、そして背景としてゲージ場である。スカラーはスピン0を持ち、1成分場\(\phi(x)\)で表される。この講義が2010年に録音されたとき、ヒッグスはまだ発見されておらず、素粒子物理学で期待される主要なスカラーとして記述されていた。その歴史的状況が「``ヒッグリー・ウィグリー''」という遊び心のある表現を説明する。その場は物理的に緊要であったが、実験的にはまだ確立されていなかった。

実スカラー場に対して、伝播関数は模式的に
\begin{equation}
G_\phi(x,y)=\expect{0|T\!\left\{\phi(x)\phi(y)\right\}|0}
=G_\phi(\Delta).
\end{equation}
と書ける。時間順序の規約と規格化はパワーカウントには重要でない。重要な事実は、2つの場の挿入が一方の点での生成と他方の点での消滅を記述することである。

\begin{classroomqa}
\audiencequestion{\(x\)と\(y\)は空間内の位置ですか、それとも時空内の位置ですか？\lecturetimestamp{00:09:31}}
\lecturerresponse{それらは時空点である。それぞれは4元ベクトルである。したがって伝播関数はそれらの4次元的な隔たりに依存し、ローレンツ不変な真空では対応する不変間隔に依存する。}
\end{classroomqa}

スピン1/2の粒子はディラック場\(\psi_i(x)\)によって記述される。添字\(i\)はスピノル添字であり、空間成分ではない。したがって伝播関数は行列である：
\begin{equation}
S_{ij}(x,y)=\expect{0|T\!\left\{\psi_i(x)\bar\psi_j(y)\right\}|0}
=S_{ij}(\Delta).
\end{equation}

聴衆への問いが成分数の答えを引き出す。4次元ディラック場は4つの複素成分をもつ。質量のないワイル場は2つの複素成分で記述できるので、左巻き場と右巻き場を別々に扱うことができる。ここでの次元の議論には、スピノル行列構造が存在することだけが重要である。

\(S_{ij}\)は\(4\times4\)行列であるから、ディラック行列の基底で展開できる。まずそれらの添字を省略し、必要になったときに復活させる。これは行列構造が存在しないという主張ではなく、紫外べき則を確立する間の意図的な単純化である。

\begin{classroomqa}
\audiencequestion{伝播関数に共役場すなわちダガー付き場はいつ現れなければならないか。\lecturetimestamp{00:14:06}}
\lecturerresponse{この概略的な議論における実の中性スカラー場では、同一の場が生成演算子と消滅演算子を含むので、別個のダガー付き場は必要ない。電荷をもつ場では、粒子と反粒子の演算子を区別しなければならない。したがって電子に対して、\(\psi\)、\(\psi^\dagger\)、および共変的に\(\bar\psi=\psi^\dagger\gamma^0\)は異なる役割を果たす。}
\end{classroomqa}

物理的なヒッグス粒子は中性であるが、その背後にある標準模型のヒッグス場は電弱二重項である。この講義では実の中性スカラー場を局所模型として用いている。この単純化は紫外の評価には十分である。\editorialnote{この区別により、概略的な実スカラーの例が電弱ヒッグス多重項についての完全な記述と誤解されることが防がれる。}

\section{自然単位とスカラー場の次元}

\(\hbar=c=1\)と置く。すると長さと時間は同じ単位をもち、エネルギー、質量、運動量は逆長さの単位をもつ：
\begin{equation}
[m]=[E]=[p]=L^{-1}.
\end{equation}
作用は位相\(e^{iS}\)に現れるので無次元である。4次元時空における代表的なスカラー作用は
\begin{equation}
S_\phi=\int d^4x\left[
\frac12\partial_\mu\phi\,\partial^\mu\phi
-\frac12m^2\phi^2
-\frac{\lambda}{4!}\phi^4
+\cdots\right].
\end{equation}
である。\(1/2\)や\(1/4!\)といった慣習的な因子は次元解析に影響しない。

運動項から始める。\([d^4x]=L^4\)および\([\partial_\mu]=L^{-1}\)であるから、無次元性より
\begin{align}
1
&=L^4\left(L^{-1}[\phi]\right)^2,\\
[\phi]&=L^{-1}.
\end{align}
が得られる。質量項は検算を与える：
\begin{equation}
L^4[m^2][\phi^2]
=L^4L^{-2}L^{-2}=1.
\end{equation}

次に会場に投げかける問いは、四次結合定数の次元である。それは四次元時空では無次元である：
\begin{equation}
L^4[\lambda][\phi]^4=L^4[\lambda]L^{-4}=1
\quad\Longrightarrow\quad [\lambda]=L^0.
\end{equation}
このことが、\(\phi^4\)相互作用を四次元において特別なものにしている。\(\phi^6\)およびそれ以上の相互作用の次元も同じ規則に従うが、それらの結合定数は質量の負の冪を伴うため、より制限された有効理論としての役割を告げている。

このように、場の次元はどの相互作用が無次元係数とともに現れうるかをすでに教えてくれる。同じ方法によって、詳細なファインマン積分を評価することなく、伝播関数と湯川結合定数の両方が決定されるであろう。

\section{短距離におけるスカラー伝播関数}

まず質量のないスカラーを考える。固有の長さスケールは存在しないので、利用できる唯一のスケールは隔たり\(\Delta\)である。並進不変性により\(\Delta\)のみが残り、ローレンツ不変性により依存性はその不変な大きさに還元され、2つのスカラー場の積は次元\(L^{-2}\)をもつ。数因子を除いて
\begin{equation}
G_0(\Delta)\sim\frac{1}{\Delta^2}.
\end{equation}
である。\(4\pi^2\)のような因子は次元の議論だけでは決められない。

質量はコンプトン長\(m^{-1}\)を導入する。その効果は分散関係
\begin{equation}
E=\sqrt{\bm p^{,2}+m^2},
\qquad
E=|\bm p|\quad(m=0).
\end{equation}
から最も容易に理解される。\(|\bm p|\gg m\)では質量は無視できる。高い運動量は短い距離を探るので、質量をもつ伝播関数は\(\Delta\to0\)のとき質量のない\(1/\Delta^2\)の振る舞いに近づかなければならない。\(m^{-1}\)よりずっと大きな距離では、伝播は抑制される。この講義はその事実を次のように概略的に表す：
\begin{equation}
G_m(\Delta)\sim \frac{e^{-m|\Delta|}}{\Delta^2}.
\end{equation}
これは意図的に厳密な式ではない。4次元ユークリッド空間では、厳密な自由伝播関数は\(mK_1(mr)/r\)に比例し、\(r=|\Delta|\)である。それは同じ\(1/r^2\)の紫外特異性をもち、べき則の前因子を伴う指数関数的な長距離減衰をもつ。\editorialnote{ベッセル関数の形は、この講義の概略的な指数関数に対する標準的な厳密化を与えるものであり、そのべき数え上げの議論には必要ない。}

決定的な教訓は単純である。質量は長距離の伝播を和らげるが、同一点の特異性を修復しない。点に十分近ければ、あらゆる有限の質量は無関係になる。

\section{スカラーのタッドポールと微調整}

四次相互作用\(\lambda\phi^4\)は四本の線が集まる頂点を生じる。その脚のうち二本を互いに結び、二本の外部スカラー脚を残す。得られるタッドポール図は\(\phi^2\)の係数、したがってスカラー質量の二乗を補正する。

保存されたフレームは、次元解析からこのループ補正への正確な移行を捉えている。分散関係は左側に残り、スカラー作用と場の次元は右側に部分的に見えたままであり、小さなタドポールがちょうど中央に現れたところである。

\begin{figure}[H]
\centering
\includegraphics[width=0.72\textwidth]{lecture_03_figure_02.png}
\lectureframe{00:28:30.560}

\medskip
\begin{tikzpicture}[line width=0.8pt,scale=1.0]
  \draw[dashed] (-2.6,0) -- (0,0) -- (2.6,0);
  \fill (0,0) circle (0.055);
  \draw (0,0) .. controls (-0.85,0.7) and (-0.85,1.6) .. (0,1.6)
                .. controls (0.85,1.6) and (0.85,0.7) .. (0,0);
\end{tikzpicture}
\caption{スカラー自己エネルギーの移行における黒板、続いてタドポール位相の保守的な概略再描画。スクリーンショットはループ、相対論的かつ質量ゼロのエネルギー公式、および部分的な作用を確実に示している。隠された係数やラベルは可視であるとは主張しない。}
\end{figure}

頂点では、スカラーは吸収されて再放出される一方、他の二本の脚はその点からそれ自身へと伝播する。数値的因子と組み合わせ的因子を除けば、補正は
\begin{equation}
\delta m_\phi^2\sim \lambda G_\phi(0).
\end{equation}
である。一致する伝播関数は連続体近似において無限大である。有効理論は任意に短い距離での妥当性を主張しないので、\(|\Delta|\geq\delta\)を課す。すると
\begin{equation}
G_\phi(0)\longrightarrow G_\phi(\delta)sim\frac{1}{\delta^2},
\qquad
\delta m_\phi^2\sim\frac{\lambda}{\delta^2}sim\lambda\Lambda^2.
\end{equation}
となる。この発散が二次と呼ばれるのは、運動量カットオフの二乗を含むからである。運動量空間での積分も同じ結論を与える。過剰な寄与は任意に大きなループ運動量から来る。最短スケールが最終的に大統一、量子重力、あるいはまだ想像されていない何かによって設定されるかどうかは、有効理論の診断を変えない。

これは微調整問題である。もし\(\delta\)が極めて小さく、\(\lambda\)が異常に小さくなければ、輻射補正は低エネルギーで観測されるスカラー質量の二乗よりもはるかに大きい。その場合、裸の係数とループ補正は信じられないほど多くの桁で相殺しなければならない：
\begin{equation}
m_{\mathrm{physical}}^2
=m_{\mathrm{bare}}^2+\delta m_\phi^2.
\end{equation}
質量を持つ伝播関数はこの状況を救わない。発散の原因となる距離では、\(e^{-m|\Delta|}\simeq1\)である。紫外ループはカットオフに到達するずっと前に質量を忘れている。

\section{フェルミオンの次元とその伝播関数}

次にディラックフェルミオンについて同じ次元解析を繰り返す。その自由作用は
\begin{equation}
S_\psi=\int d^4x\,
\bar\psi\left(i\gamma^\mu\partial_\mu-m\right)\psi.
\end{equation}
である。ガンマ行列は数値的かつ無次元である。運動項は一つの微分しか含まないので、
\begin{align}
1
&=L^4[\bar\psi]L^{-1}[\psi],\\
[\psi]=[\bar\psi]&=L^{-3/2}.
\end{align}
となる。半整数の冪は驚くべきものではない。フェルミオンは多くの場所に半整数を持ち込む。

\begin{classroomqa}
\audiencequestion{質量項から\([\psi]\)を導出する方が簡単ではないか？\lecturetimestamp{00:36:02}}
\lecturerresponse{はい。質量項は\(L^4L^{-1}[\bar\psi\psi]=1\)を与え、したがって\([\bar\psi\psi]=L^{-3}\)および\([\psi]=L^{-3/2}\)となる。運動項からの導出を先に用いたのは、それが質量ゼロの理論でも意味を持ち続け、フェルミオン質量が\(m^2\)ではなく\(m\)として現れる理由を示すからである。}
\end{classroomqa}

二つのフェルミオン場は次元\(L^{-3}\)を持つので、質量ゼロの伝播関数に対する講義の簡潔な次元の略記は
\begin{equation}
S(\Delta)\sim\frac{1}{|\Delta|^3}.
\end{equation}
である。その行列構造を復元すると、ローレンツ共変性が要求するより正確な分子が得られる：
\begin{equation}
S_{ij}(\Delta)
\propto
\frac{\left(\gamma_\mu\Delta^\mu\right)_{ij}}
{\left(\Delta^2\right)^2}.
\end{equation}
分子は長さの冪を一つ加え、したがって分母は間隔の四乗を含む。正味の次元は\(L^{-3}\)のままである。スカラーの場合と同様に、フェルミオン質量は大きな分離で伝播関数を抑制するが、この短距離スケーリングは変えない。

\section{閉じたフェルミオンループが負である理由}

ここで、フェルミオン的ダイアグラムが悪いボソン的自己エネルギーを相殺するかどうかを問うことができる。スカラーのタドポールは、\(\pi\)の冪や組み合わせ係数のような正の数値因子を除いて、正の符号を持つ。閉じたフェルミオンループは反対の符号を持つ。

\begin{classroomqa}
\audiencequestion{フェルミオンループのマイナス符号は、以前の一回転の議論に関係しているのか？\lecturetimestamp{00:40:21}}
\lecturerresponse{それは同じフェルミオン的符号構造に属するが、ここで有用な証明はループの幾何学的向きではなく交換反対称性を通じて進む。二つのループを切り、同一粒子を交換し、それらを再接続する。交換は二つのループを一つに変え、フェルミオンに対してのみマイナス符号を与える。}
\end{classroomqa}

この議論は完全に述べる価値がある。一ループダイアグラムが符号\(\sigma\)を持つと仮定する。プラスかマイナスかのいずれかである。二つの非連結な同一ループは符号
\begin{equation}
\sigma^2=+1.
\end{equation}
を持つ。両方のループを同時に切る。ここで二つの同一粒子が切断の間を伝播している。それらの通常の二粒子波動関数は、ボソンに対して対称であり、フェルミオンに対して反対称である：
\begin{align}
\Psi_B(x_1,x_2)&=+\Psi_B(x_2,x_1),\\
\Psi_F(x_1,x_2)&=-\Psi_F(x_2,x_1).
\end{align}
反対称性はパウリ原理も説明する。もし両方のフェルミオンが同じ状態に置かれるなら、交換は物理的に何も変えないが、波動関数を反転させなければならず、それが消えることを強制する。

二つの粒子を交換した後、切断された線を再接続する。位相幾何学的には、二つのループは一つになった。交換しなかった二ループ振幅は正であった。交換により、ボソンでは正のままであり、フェルミオンでは符号が反転する。したがって
\begin{equation}
\begin{aligned}
\operatorname{sign}(\text{closed boson loop})&=+1,\\
\operatorname{sign}(\text{closed fermion loop})&=-1.
\end{aligned}
\end{equation}
この結論は一般的である。各々の独立な閉じたフェルミオンループは\(-1\)の因子を与える。

\section{湯川ループ}

スカラー質量に戻る。ボソンは単一のフェルミオンに転換できない。フェルミオン数パリティが変化するからである。最も単純なフェルミオン補正では、スカラーがフェルミオン・反フェルミオン対に分裂し、それがループを伝播して再結合する。相互作用は湯川結合であり、
\begin{equation}
\mathcal L_Y=-g\phi\bar\psi\psi.
\end{equation}
ヒッグスの場合、フェルミオンはクォークでもレプトンでもよい。相互作用点は二つあり、したがって\(g\)の因子も二つある。

スカラーのタドポールは一点を含み、その全体の位置は固定できた。並進不変性は時空体積を与えるだけだからである。湯川ループは二点を含む。一方の点を固定してもよく、あるいは等価にそれらの平均位置を固定してもよいが、それらの相対的な分離\(\Delta\)について積分しなければならない。各フェルミオン伝播関数は短距離因子\(1/|\Delta|^3\)に寄与するので、
\begin{equation}
\delta m_{\phi,F}^2
\sim
-g^2\int_{|\Delta|\geq\delta}d^4\Delta\,
\frac{1}{|\Delta|^6}.
\end{equation}

二つの伝播関数、分離の積分、そして二つの結合が書かれた後、最後の因子が場から求められる。答えはマイナス符号である。二つの伝播関数は閉じてフェルミオンループになるので、それ以外では正である短距離積分全体に\(-1\)が掛かる。

次元解析がすでに発散を決定している。長さの四乗は\(d^4\Delta\)から来て、六乗は分母に現れ、残る唯一の長さのスケールはカットオフである：
\begin{equation}
\int_{|\Delta|\geq\delta}
\frac{d^4\Delta}{|\Delta|^6}
\sim\frac{1}{\delta^2}.
\end{equation}
四次元ユークリッド動径座標では、これを明示的にできる：
\begin{equation}
\int_{\delta}^{R}d^4\Delta\,\frac{1}{|\Delta|^6}
=2\pi^2\int_{\delta}^{R}dr\,\frac{r^3}{r^6}
=\pi^2\left(\frac{1}{\delta^2}-\frac{1}{R^2}\right).
\end{equation}
スピノルのトレースと伝播関数の規格化は係数を変えるが、カットオフの二次の冪は変えない。\editorialnote{動径積分は、講義で行われた次元評価の標準的な明示版である。それはユークリッド距離で書かれており、そこでは正値性と紫外の冪 counting が透明である。}

\begin{classroomqa}
\audiencequestion{積分の極限を逆にしたり選んだりすると符号が変わりうるか？ \lecturetimestamp{00:51:06}}
\lecturerresponse{いいえ。これは極限が物理的な符号を生み出す一次元の向き付けられた積分ではない。これは四つの成分にわたる積分、あるいはより明快には正の動径測度にわたる積分である。その紫外の大きさは正である。負の寄与は閉じたフェルミオンループからのみ来る。}
\end{classroomqa}

\begin{classroomqa}
\audiencequestion{湯川結合\(g\)の次元は何か？ \lecturetimestamp{00:51:51}}
\lecturerresponse{四次元では無次元である。\([\phi]=L^{-1}\)かつ\([\bar\psi\psi]=L^{-3}\)なので、積\(\phi\bar\psi\psi\)は次元\(L^{-4}\)を持ち、\(d^4x\)によって正確に相殺される。}
\end{classroomqa}

したがって主要なフェルミオン補正は
\begin{equation}
\delta m_{\phi,F}^2
\sim -C_F\frac{g^2}{\delta^2}
=-C_Fg^2\Lambda^2,
\end{equation}
ここで\(C_F>0\)は数値的因子、スピン因子、多重度因子を含む。スカラーループは対応する形
\begin{equation}
\delta m_{\phi,B}^2
\sim+C_B\frac{\lambda}{\delta^2}
=+C_B\lambda\Lambda^2.
\end{equation}
を持つ。我々はまだ微調整問題を解決していないが、不可欠なものを見つけた。同じ紫外次数と同じ反対符号を持つ寄与である。

\section{相殺が要求するもの}

無関係な結合が、スケールの大きな階層によって要求される精度に正確に従う
\begin{equation}
C_B\lambda=C_Fg^2
\end{equation}
というのは、並外れた偶然であろう。講義ではこれを\(g^2=\lambda\)と略記し、規約に依存する因子と多重度因子を抑制している。真の解決には、関係を単に許すのではなく強制する数学的構造が必要である。

それが超対称性の動機である。超対称性はボソンとフェルミオンを関係づけ、それらの結合を、多くの危険な発散項を相殺させる組み合わせに制限する。相殺は、パラメータが一ループずつ調整されたからではなく、対称性から従うために、図ごとに繰り返される。

質量の一致も重要である。ボソンとフェルミオンの質量が同一であれば、それらの完全な伝播関数を最も緊密に対にすることができる。質量がわずかに異なると、完全な有限の答えは厳密には相殺しない。なぜなら、伝播関数はそれらのコンプトン波長と同程度の距離で異なるからである。それでも、それらに共通する短距離漸近形は発散部分を相殺しうる：
\begin{equation}
\mathcal I_B(\Delta)-\mathcal I_F(\Delta)
\quad\hbox{may be finite as}\quad \Delta\to0,
\end{equation}
ここで\(\mathcal I_B\)と\(\mathcal I_F\)は、頂点、分子、および多重度を含めた後の完全なボソン的およびフェルミオン的ループ被積分関数を表す。伝播関数自体は異なる次元を持ち、直接差し引かれているのではない。紫外項が相殺しても、有限の残差が残る場合がある。これが近似された、あるいは破れた超対称性への入り口である。厳密な等号は最もきれいな相殺を与える；十分に制御された破れであれば、観測可能な質量分裂を残しつつ紫外発散を除去できる場合がある。

したがって、我々はその代数をまだ書かずに超対称性に到達した。推論の連鎖は完結している：
\begin{equation}
\begin{gathered}
G_\phi(\Delta)\sim \Delta^{-2},
\qquad
\delta m_{\phi,B}^2\sim+C_B\lambda\Lambda^2,\\[0.4em]
S_\psi(\Delta)\sim \Delta^{-3},
\qquad
\delta m_{\phi,F}^2\sim-C_Fg^2\Lambda^2,\\[0.4em]
\text{symmetry-enforced matching}\\
\Longrightarrow\quad
\text{quadratic divergences cancel}.
\end{gathered}
\end{equation}
次の課題は、それらの関係を強制する抽象的な対称性原理を理解し、それがどのように作用するかの具体例を見ることである。